% \section{Related Work}
% \paragraph{LLM/VLM-Powered GUI Agents.}
% The recent proliferation of LLMs and VLMs has spurred growing interest in developing GUI agents capable of performing challenging tasks across digital devices, such as mobile phones~\citep{rawles2024androidworld}, web applications~\citep{zhou2023webarena, deng2023mind2web}, and desktops~\citep{xie2024osworld}.
% Early research focused on designing autonomous agent frameworks~\citep{zheng2024gpt, rawles2024androidworld, jia2024agentstore} that prompt commercial models to interact with operating systems via code generation~\citep{sun2024survey} or tool use~\citep{zhang2025appagent}.
% UFO~\citep{zhang2024ufo} decomposes user instructions and generates executable Python calls to perform tasks on Windows systems.
% OS-Copilot~\citep{wu2024copilot} introduces a tool memory mechanism that facilitates agent self-improvement.

% With rising demand for open-source and customizable agents, a parallel line of work aims to train LLMs/VLMs to enhance their agentic capabilities—such as understanding GUI environments, planning and self-reflection~\citep{xu2024aguvis, qin2025ui}, and action execution~\cite{chen2024guicourse, chai2024amex}.
% Substantial efforts have been devoted to collecting GUI-centric training corpora to enrich model knowledge, including OCR~\citep{lu2024gui}, interface summarization~\citep{you2024ferret}, and question answering on web pages~\citep{chen2024guicourse}, and, crucially, GUI task trajectories for training versatile GUI agents~\citep{deng2023mind2web, rawles2023androidinthewild, zhang2024android, zheng2025vem, sun2024genesis, zhang2025breaking}.
% A central requirement of agent development is the ability to interact with realistic GUI environments deployed in virtual machines and Chrome-based browsers.
% Early explorations relied on structured environment metadata—executing actions by selecting elements from HTML or accessibility trees~\citep{zhou2023webarena, deng2023mind2web, he2024webvoyager, zhang2025appagent}. However, this approach is inefficient and difficult to generalize across diverse operating systems and applications~\citep{cheng2024seeclick, xu2024aguvis}.
% Consequently, recent efforts have shifted toward a unified vision-based paradigm, in which agents interact with GUIs by purely perceiving the screen and executing mouse and keyboard actions~\citep{openai2025operator, anthropic2024cuda}, much like humans do.
% A key challenge in this paradigm is mapping natural language instructions to on-screen GUI elements for click actions, which we refer to as \textit{Visual Grounding for GUI Agents}.

% \paragraph{Visual Grounding for GUI Agents.}
% Given a screenshot and a natural language instruction or description, GUI visual grounding aims to identify the location of the corresponding element on the screen.
% While analogous to visual grounding in natural images, this task poses new challenges, as conventional methods struggle in GUI contexts due to the substantial gap between screenshots and natural images~\citep{cheng2024seeclick, gou2024navigating}.
% One solution is to fine-tune object detection models on GUI data. For example, OmniParser~\citep{lu2024omniparser} parses interface screenshots into DOM-like structures, enabling seamless integration with SoM-style grounding frameworks~\citep{yang2023set}.

% A more widely adopted approach frames GUI visual grounding as a text-based coordinate generation task, enabling VLMs to directly output the target location in natural language format~\citep{chen2023shikra, wang2024qwen2}.
% Following the GUI grounding pre-training strategy~\citep{cheng2024seeclick, hong2024cogagent}, extensive efforts have explored scaling data and models, leading to significant performance gains~\citep{lin2024showui, yang2024aria, pawlowski2024tinyclick, tang2025think, bai2025qwen2, qin2025ui}.
% UGround~\citep{gou2024navigating} introduces a data pipeline for synthesizing diverse and realistic grounding data from webpages.
% OS-Atlas~\citep{wu2024atlas} develops a multi-platform data infrastructure and builds a foundation action model that unifies the action space for GUI agents.
% Moreover, \citep{xu2025attention} proposes a training-free approach that performs GUI grounding by leveraging the internal attention of VLMs.
% In this paper, we propose a compelling alternative to prevailing coordinate-centric methods. Our approach introduces an action token that directly attends to the corresponding image patches, enabling coordinate-free GUI visual grounding.


\section{Related Work}

\paragraph{LLM/VLM-Powered GUI Agents.}
The advent of LLMs and VLMs has catalyzed the development of GUI agents that can understand natural language instructions and perform complex tasks across mobile~\citep{rawles2024androidworld}, web~\citep{zhou2023webarena, deng2023mind2web}, and desktop environments~\citep{xie2024osworld, zhang2024ufo, zhang2025ufo2}.
% Early efforts focused on autonomous frameworks that prompt commercial models to execute tasks via code generation~\citep{sun2024survey} or external tool use~\citep{zhang2025appagent}, such as GPT-Engineer~\citep{zheng2024gpt}, AgentStore~\citep{jia2024agentstore}, and AndroidWorld~\citep{rawles2024androidworld}. UFO~\citep{zhang2024ufo} enables structured task decomposition and Python-based execution on Windows, while OS-Copilot~\citep{wu2024copilot} introduces tool memory for iterative refinement.
Early research focused on designing autonomous agent frameworks~\citep{zheng2024gpt, wu2024copilot, jia2024agentstore} that prompt commercial models to interact with operating systems via code generation~\citep{sun2024survey, sun2025scienceboard} or tool use~\citep{zhang2025appagent,zhang2025api}.
With rising demand for open-source and customizable agents, a parallel line of work focuses on training LLMs/VLMs for enhanced agentic capabilities, including GUI understanding, planning, and execution~\citep{xu2024aguvis, qin2025ui, chen2024guicourse, chai2024amex}. The key to these efforts is collecting GUI-specific training data, such as OCR annotations~\citep{lu2024gui}, interface summaries~\citep{you2024ferret}, QA pairs~\citep{chen2024guicourse}, and large-scale task demonstrations~\citep{deng2023mind2web, rawles2023androidinthewild, zhang2024android, zheng2025vem, sun2024genesis, cheng2024vision, zhang2025breaking}.
% These datasets often require executing actions within realistic environments, typically powered by virtual machines or Chrome-based automation.

A central requirement of agent development is the ability to interact with realistic GUI environments deployed in virtual machines and Chrome-based browsers.
While early agents operated over structured metadata like HTML or accessibility trees~\citep{zhou2023webarena, he2024webvoyager}, such representations are brittle and inconsistent across platforms~\citep{cheng2024seeclick, xu2024aguvis}. Consequently, recent trends have shifted toward a \emph{vision-centric paradigm}, where agents interact with raw screenshots using mouse and keyboard inputs~\citep{openai2025operator, anthropic2024cuda}, closely mimicking human behavior. Within this setting, a central challenge emerges: grounding natural language instructions to specific GUI regions, referring to as \textit{GUI Visual Grounding}.

\paragraph{GUI Visual Grounding.}
Given a GUI screenshot and a natural language instruction, GUI visual grounding aims to locate the target region for interaction.
Although conceptually related to grounding in natural images, this task presents unique challenges due to the semantic density and structural regularity of GUI layouts~\citep{cheng2024seeclick, gou2024navigating}. A common approach frames GUI grounding as a text-based coordinate prediction task, where models generate point positions (\eg, \texttt{x=..., y=...}) as output language tokens~\citep{chen2023shikra, wang2024qwen2}. This formulation has led to widespread adoption due to its simplicity and compatibility with existing LLMs/VLMs. To improve performance, prior works have scaled both models and training data~\citep{lin2024showui, yang2024aria, pawlowski2024tinyclick, tang2025think, bai2025qwen2, qin2025ui, xie2025scaling}. UGround~\citep{gou2024navigating} proposes a data pipeline for synthesizing diverse GUI grounding examples, while OS-Atlas~\citep{wu2024atlas} offers a multi-platform dataset and a unified GUI action model. More recently, \citet{xu2025attention} introduced a training-free approach that performs GUI grounding by leveraging the internal attention of VLMs.

Despite their success, coordinate-based methods suffer from key limitations,
including weak spatial inductive bias, ambiguous point supervision, and resolution mismatches between visual features and action targets.
% they lack spatial inductive bias, struggle with ambiguous supervision, and suffer from resolution mismatches between visual features and action targets.
This paper presents a compelling alternative to the prevailing coordinate-based method: GUI-Actor, a novel \textit{coordinate-free grounding framework} for GUI agents.
It introduces an \texttt{<ACTOR>} token that attends directly to relevant image patches via an attention-based action head, enabling more human-like grounding while mitigating the limitation of coordinate-based methods.
% In contrast, our work proposes a novel, \textit{coordinate-free grounding approach} by introducing an \texttt{[ACTOR]} token that attends directly to relevant image patches via an attention-based action head.
% This design enables more human-like grounding behavior, improves spatial-semantic alignment, and handles ambiguous supervision more gracefully.

% Some works attempt to mitigate this by predicting bounding boxes~\citep{lu2024omniparser}, but they still rely on raw coordinate outputs without spatially grounded representations.



% A grounding verifier further refines predictions by selecting the most plausible region, resulting in more robust and generalizable GUI action execution.
